% Full provenance: HIBRIDACION_ACCION_MASA_PROPAGACION.md, §§2--7.
% Editorial records, history and source locators are preserved in gestion/.
\section{Action, mass and propagation: structural compatibility and transport}
\label{md:hib:seccion}

\subsection{Shared action section and mass law}
\label{md:hib:accion-masas}

Let
\begin{equation}
 s_0=10^{-34}\mathcal U_S,\qquad
 \hbar_{\mathrm{ret}}=\eta_{\mathrm{ret}}s_0,\qquad
 R_{\mathrm{act}}=\frac{H_5}{\eta_{\mathrm{ret}}}.
 \label{md:hib:accion}
\end{equation}
The coordinates constructed in degrees are retained:
\begin{equation}
 A_{\deg}=1000\alpha,\qquad
 C^*_{\deg}=2(\eta_{\mathrm{ret}}+\alpha).
 \label{md:hib:angulos}
\end{equation}
Only afterwards are they transported to radians through
$x=\pi A_{\deg}/180$ and $y=\pi C^*_{\deg}/180$. The channels are
$q_\pm=\exp(-x\mp y)$, with $x>|y|$. The elliptic shape is expressed by
\begin{equation}
 \zeta=\operatorname{artanh}(y/x).
 \label{md:hib:forma}
\end{equation}
The composition first uses the dimensionless coefficient
$\eta_{\mathrm{ret}}$, thereby preserving the types of action and angle.
In this notation, $\eta_{\mathrm{ret}}$ is the action coefficient,
$\eta_h(\Gamma)$ belongs to a route tower, and $\zeta$ expresses the
deformation of the ellipse; their roles remain distinct.

The working domain has $t_0,c_{\mathrm{int}},\eta_{\mathrm{ret}}>0$, a
convergent character and a factor $\Xi_\Gamma>0$. The constructed mass
readout operator determines
\begin{equation}
 \omega_0=\frac{2\pi}{108t_0},\qquad
 E_{\mathrm{clk}}=\hbar_{\mathrm{ret}}\omega_0,\qquad
 m_{\mathrm{clk}}=\frac{E_{\mathrm{clk}}}{c_{\mathrm{int}}^2}.
 \label{md:hib:reloj}
\end{equation}
For a route with a scalar realization,
\begin{equation}
 m_\Gamma^\beta=m_{\mathrm{clk}}\Xi_\Gamma,\qquad
 \Xi_\Gamma=b_e\,
 \chi_{\mathrm{ref}}(\sigma_\Gamma-\sigma_e,\eta_\Gamma-\eta_e)
 R_{\mathrm{act}}^{\kappa_\Gamma}.
 \label{md:hib:factor-ruta}
\end{equation}
$\Xi_\Gamma$ abbreviates the structural route factor. The following
normalization is explicitly retained:
\begin{equation}
 b_e=\frac{\varepsilon_e^{(0)}\mathcal U_E}{E_{\mathrm{clk}}}.
 \label{md:hib:normalizacion-electron}
\end{equation}
The central state retains its own normalization relative to the
chronological quantum:
\begin{equation}
 \varepsilon_e^{(0)}=\frac{\sqrt3}{4}
 \left[\exp\!\left(\frac{\varphi}{\pi^2}\right)-22\alpha^3\right]
 (1+15\Delta_4),\qquad
 \varepsilon_e^\beta=\varepsilon_e^{(0)}
 R_{\mathrm{act}}^{80-54\alpha+6\Delta_4}.
 \label{md:hib:electron}
\end{equation}
The refined character includes the terms
\begin{equation}
 n_Ax+\frac{s_C}{6}y+k_{\Delta_4}\Delta_4+b_\zeta\zeta_{270}
 +\sum_h N_h(q_-^h-q_+^h),\qquad
 N_h=\eta_h+b_{120}\mathbf 1_{\{h=120\}},
 \label{md:hib:caracter}
\end{equation}
with $h$ in the semigroup generated by $90$ and $120$. The height-$120$
term is included once. The primitive return coefficients are constructed
through
\begin{equation}
 \begin{aligned}
 a_n(\Gamma)&=\operatorname{Tr}(R_\Gamma U_\Gamma^n),\\
 p_n(\Gamma)&=\frac1n\sum_{d\mid n}\mu_{\mathrm{Mob}}(d)
                         a_{n/d}(\Gamma),\\
 \eta_h(\Gamma)&=p_h(\Gamma)-p_h(\Gamma_e),
 \end{aligned}
 \label{md:hib:retornos-primitivos}
\end{equation}
where $\mu_{\mathrm{Mob}}$ is the arithmetic Möbius function.
Action participates in the absolute scale, in the ratio $R_{\mathrm{act}}$
and in $y$, which determines the channels. Consequently, cancellation of
a dimensional scale retains the remaining structural dependencies.
When the clock varies, $b_e=\varepsilon_e^{(0)}\mathcal U_E/E_{\mathrm{clk}}$
varies with it: preserving the electronic readout requires transporting
the complete character, rather than fixing that quotient separately.
The integrated longitudinal realization uses
$t_0=\pi_{\rm HMT}L_\alpha/(54c_{\mathrm{int}})$ from
\eqref{md:rigidez:reloj}.

\subsubsection{Frequency and length of each massive section}
\label{md:hib:frecuencia-longitud}

Define $E_\Gamma=m_\Gamma c_{\mathrm{int}}^2$ as rest energy,
$\Omega_\Gamma=E_\Gamma/\hbar_{\mathrm{ret}}$,
$\nu_\Gamma=\Omega_\Gamma/(2\pi)$ and
$R_{108}=c_{\mathrm{int}}/\omega_0$. Then
\begin{equation}
 \begin{aligned}
 E_\Gamma&=\hbar_{\mathrm{ret}}\omega_0\Xi_\Gamma,&
 \Omega_\Gamma&=\omega_0\Xi_\Gamma,\\
 \nu_\Gamma&=\frac{\Xi_\Gamma}{108t_0},&
 \bar\lambda_\Gamma&=
 \frac{\hbar_{\mathrm{ret}}}{m_\Gamma c_{\mathrm{int}}}
 =\frac{R_{108}}{\Xi_\Gamma}.
 \end{aligned}
 \label{md:hib:frecuencias}
\end{equation}
\begin{proof}
Substitute $m_\Gamma=m_{\mathrm{clk}}\Xi_\Gamma$ and perform the
cancellations in the same dimensional chart. For positive routes $X,Y$,
the character property and common normalization give
\begin{equation}
 \frac{\Omega_Y}{\Omega_X}=\frac{m_Y}{m_X}
 =\chi_{\mathrm{ref}}(\sigma_Y-\sigma_X,\eta_Y-\eta_X)
 R_{\mathrm{act}}^{\kappa_Y-\kappa_X}.
 \label{md:hib:cocientes}
\end{equation}
\end{proof}
The cancellation is dimensional; the channels and relative memory remain.
The zero sector lies outside the domain of inverse mass.
In a nonscalar fibre, the operator or its spectral measure is retained,
rather than selecting an arbitrary point representative.

\subsubsection{Lifted frequency and compatibility of normalizations}
\label{md:hib:tres-frecuencias}

The chronological construction also produces the lifted holonomy frequency
\begin{equation}
 \omega_j=\frac{\theta_j+2\pi w_j}{108t_0},\qquad
 \lambda_j=r_j\exp(i\theta_j),
 \label{md:hib:frecuencia-levantada}
\end{equation}
where $w_j$ retains the winding of the history. Its energy readout is
\begin{equation}
 E_j=\hbar_{\mathrm{ret}}\omega_j
 \chi_{\mathrm{full}}R_{\mathrm{act}}^{\kappa_j}.
 \label{md:hib:energia-levantada}
\end{equation}
This frequency, $\omega_0$ and $\Omega_\Gamma$ correspond to distinct
readout operators. If both energy expressions represent the same mode
and the same exponent, their exact compatibility requires
\begin{equation}
 \frac{\omega_j}{\omega_0}\chi_{\mathrm{full}}
 =b_e\chi_{\mathrm{ref,rel}}.
 \label{md:hib:compatibilidad-normalizaciones}
\end{equation}
The two normalizations are identified through this equality and are not
multiplied together. The reduced phase $\theta_j$ alone does not recover
$w_j$. A phase return therefore remains distinct from the complete
chronological history. Winding retains its construction map; this
distinction does not authorize choosing it freely or producing energy
without a dynamics.

\subsection{Spatial reciprocity and gravitational composition}
\label{md:hib:gravedad}

The circular gravitational section expresses the result of
Theorem~\ref{md:rigidez:teorema}:
$G_a=c_{\mathrm{int}}^3L_\alpha^2/\hbar_a$ and $R_{108}=L_\alpha$.
The length arises from the inscribed return and precedes this
circular identification. Retaining
$\hbar_a=\hbar_{\mathrm{ret}}$ and the reduced convention
$r_g=Gm/c_{\mathrm{int}}^2$, the preceding mass determines
\begin{equation}
 \boxed{r_{g,\Gamma}=R_{108}\Xi_\Gamma,\qquad
 \bar\lambda_\Gamma=\frac{R_{108}}{\Xi_\Gamma},\qquad
 r_{g,\Gamma}\bar\lambda_\Gamma=R_{108}^2.}
 \label{md:hib:reciprocidad-radial}
\end{equation}
\begin{proof}
Substitution of $G_a$ and $m_\Gamma$ gives
\[
 \frac{G_am_\Gamma}{c_{\mathrm{int}}^2}
 =\frac{c_{\mathrm{int}}^3R_{108}^2}{\hbar_{\mathrm{ret}}}
 \frac{\hbar_{\mathrm{ret}}\omega_0\Xi_\Gamma}
 {c_{\mathrm{int}}^4}
 =R_{108}\Xi_\Gamma.
\]
The second equality follows from \eqref{md:hib:frecuencias}; their
product proves the third.
\end{proof}
The same route coordinate enlarges one radial readout and contracts the
other. In logarithmic form,
\begin{equation}
 \log\frac{r_{g,\Gamma}}{R_{108}}=\log\Xi_\Gamma,
 \qquad
 \log\frac{\bar\lambda_\Gamma}{R_{108}}=-\log\Xi_\Gamma.
 \label{md:hib:radios-logaritmicos}
\end{equation}
The geometric centre of the two readouts is fixed by the area $R_{108}^2$.

In the bounded operator version, let $X$ be bounded and self-adjoint, with
$X\geq\delta I$ for some $\delta>0$, on a massive fibre. The operators
\begin{equation}
 M=m_{\mathrm{clk}}X,\qquad
 \Lambda=R_{108}X^{-1},\qquad R_g=R_{108}X
 \label{md:hib:operadores-radiales}
\end{equation}
satisfy $R_g\Lambda=\Lambda R_g=R_{108}^2I$ throughout the space.
The claim uses functional calculus for the same operator. The nullspace
is excluded from the inverse. The present statement retains the bounded,
uniformly positive regime; extension to unbounded operators requires
separate proofs for the domains, compatible truncations, and passage
to the limit. Section~\ref{md:rigidez:variaciones} also proves preservation
of both identities under noncommutative variations and compatible
memory reduction.

The proposed interpretation characterizes the relative mass factor as
the reciprocal-dilation parameter of these two readouts. The circular
condition belongs to the result. This equality alone does not identify
the related lengths with measured material radii or with an electronic
horizon. The radial identity in isolation is used for its role in the
composition, without attributing independent historical priority to it.

\subsection{Constitutive response, orientation and magnitude}
\label{md:hib:constitutivas}

With $T=\exp(-xI-yR)$, $R=P_+-P_-$ and the sheet projectors retained,
the temporal readout operator is
\begin{equation}
 \begin{aligned}
 V&=(I-T^{90})(I-T^{120})^{-1}=r_+P_++r_-P_-,\\
 r_\pm&=f(s_\pm),\qquad
 s_\pm=\exp[-30(x\pm y)],\qquad
 f(s)=\frac{1-s^3}{1-s^4}.
 \end{aligned}
 \label{md:hib:operador-constitutivo}
\end{equation}
This recovers $\widehat\varepsilon=r_+^2$, $\widehat\mu=r_-^2$,
$\widehat Z=r_-/r_+$ and $\widehat c=(r_+r_-)^{-1}$.
The constitutive speed relation establishes
$c_{\mathrm{int}}=c_{\mathrm{vac}}=\widehat c\,u_C$ in the same chart.
The oriented Catalan collection participates in $f$ through its resolvent
moment, retaining the structure that precedes the endpoint value.

The exact decomposition
\begin{equation}
 V=\widehat c^{-1/2}
 \exp\!\left[-\frac12(\log\widehat Z)R\right]
 \label{md:hib:descomposicion-v}
\end{equation}
distinguishes common dilation from distribution between sheets. Sheet
exchange preserves $\widehat c$ and inverts $\widehat Z$; the Barbero
functional changes sign. Section~\ref{md:compos:difeomorfismo} proves that
$(\widehat c,\gamma)$ recovers $r_+,r_-$ in the normalized collection.
Restriction to the generated image preserves this recoverability; the
inversion statement does not declare every synthetic pair to be generated.

\subsubsection{Convexity and evaluation of the separated channels}
\label{md:hib:convexidad}

\begin{proposition}
\label{md:hib:proposicion-convexidad}
For fixed $x>0$ and $|y|<x$, the function
$y\mapsto\log\widehat c(x,y)$ is even and strictly convex.
Consequently,
\begin{equation}
 \widehat c(x,y)\geq\widehat c(x,0)
 =\frac{1}{f(e^{-30x})^2},
 \label{md:hib:cota-velocidad}
\end{equation}
with equality precisely when $y=0$.
\end{proposition}
\begin{proof}
Let $g(u)=\log f(e^{-u})$ for $u>0$. Then
\begin{equation}
 g'(u)=\frac{3}{e^{3u}-1}-\frac{4}{e^{4u}-1}>0,
 \label{md:hib:derivada-g}
\end{equation}
\begin{equation}
 g''(u)=-\frac{9}{4\sinh^2(3u/2)}
 +\frac{16}{4\sinh^2(2u)}<0.
 \label{md:hib:segunda-g}
\end{equation}
The first inequality follows because
$a/(\exp(au)-1)$ decreases in $a>0$. For the second, the function
$a/\sinh(au/2)$ is strictly decreasing: its derivative has the sign of
$\sinh v-v\cosh v$, which is negative for $v>0$. Since
\[
 \log\widehat c(x,y)=-g(30(x+y))-g(30(x-y)),
\]
its second derivative with respect to $y$ is
\begin{equation}
 -900\bigl[g''(30(x+y))+g''(30(x-y))\bigr]>0.
 \label{md:hib:segunda-velocidad}
\end{equation}
Evenness fixes the unique minimum at zero.
\end{proof}
One also obtains
\begin{equation}
 \begin{aligned}
 \log\widehat c(x,y)
 &=-2g(30x)-900g''(30x)y^2+O(y^4),\\
 \log\widehat Z(x,y)
 &=-60g'(30x)y+O(y^3).
 \end{aligned}
 \label{md:hib:desarrollos-canales}
\end{equation}
Speed loses the orientation sign but retains a quadratic response to
channel separation. Impedance retains oriented information from first
order. Averaging $x+y$ and $x-y$ before applying the readout operator
removes that contribution. The claim is a proven consequence of the
constitutive response, distinct from a new measurement.

On the positive action branch,
$\eta_{\mathrm{ret}}=(90/\pi)(y-x/500)$; hence,
\begin{equation}
 \hbar_{\mathrm{ret}}=s_0\frac{90}{\pi}(y-x/500),\qquad
 E_k=s_0\frac{90u_C|k|}{\pi}(y-x/500)\widehat c(x,y).
 \label{md:hib:energia-angular}
\end{equation}
For fixed $x$, $|k|>0$ and $x/500<y<x$, the latter function is strictly
increasing in $y$. This is a property of the homogeneous radiative
realization collection. Its application to a generated trajectory retains the
construction map for $y$; variation of a parameter of this collection is not
in itself an experimental variation of a universal constant.
Under orientation reversal, the action label is transported as well,
retaining the corresponding positive branch.

\subsection{Action, clock and thermal information}
\label{md:hib:termica}

The thermal section retains
$k_B\Theta_{108}\log3=E_{\mathrm{clk}}$, with
$\Theta_{108}=\Theta_{\mathrm{clk,ret}}$ in the return section.
Define $\beta_{108}=(k_B\Theta_{108})^{-1}$. For
$E_\Gamma=E_{\mathrm{clk}}\Xi_\Gamma$, it follows exactly that
\begin{equation}
 \beta_{108}E_\Gamma=(\log3)\Xi_\Gamma,
 \qquad
 \boxed{\exp(-\beta_{108}E_\Gamma)=3^{-\Xi_\Gamma}.}
 \label{md:hib:peso-termico}
\end{equation}
The expression is an unnormalized Gibbs weight. For a finite collection of
states, or a collection with a convergent partition sum, and multiplicities
$g_\Gamma$, one obtains
\begin{equation}
 Z_{\mathrm{part}}=\sum_\Gamma g_\Gamma3^{-\Xi_\Gamma},\qquad
 p_\Gamma=\frac{g_\Gamma3^{-\Xi_\Gamma}}{Z_{\mathrm{part}}},\qquad
 \frac{p_Y/g_Y}{p_X/g_X}=3^{-(\Xi_Y-\Xi_X)}.
 \label{md:hib:particion}
\end{equation}
A canonical realization at $\Theta_{108}$ and a common energy origin are
assumed. When a chemical potential is involved, it is incorporated into
the corresponding energy. The index $\Gamma$ enumerates distinct physical
states or groups; multiple equivalent histories do not automatically
produce a multiplicity. The dimensionless entropies are
\begin{equation}
 \begin{aligned}
 H_{\mathrm{micro}}
 &=\log Z_{\mathrm{part}}+(\log3)\sum_\Gamma p_\Gamma\Xi_\Gamma,\\
 H_{\mathrm{grupos}}
 &=H_{\mathrm{micro}}-\sum_\Gamma p_\Gamma\log g_\Gamma.
 \end{aligned}
 \label{md:hib:entropias}
\end{equation}
Finiteness of entropy additionally requires
\begin{equation}
 \sum_\Gamma g_\Gamma\Xi_\Gamma3^{-\Xi_\Gamma}<\infty.
 \label{md:hib:condicion-entropia}
\end{equation}
Finiteness of $Z_{\mathrm{part}}$ alone ensures normalization.
Both conditions hold for finite collections.

The formula links mass structure with proper frequency and thermal
weighting in the same framework. It preserves the distinction between
energy and heat, between the real coordinate $\Xi_\Gamma$ and an integer
trit count, and between a realization at $\Theta_{108}$ and a claim about
the temperature of the universe. For free bosonic or fermionic occupations
per mode, with zero chemical potential, the same substitution gives
\begin{equation}
 \overline n_\Gamma=\frac{1}{3^{\Xi_\Gamma}\mp1}.
 \label{md:hib:ocupaciones}
\end{equation}
The minus sign corresponds to Bose and the plus sign to Fermi. The total
for the degenerate group is $g_\Gamma\overline n_\Gamma$. For Bose,
$\Xi_\Gamma>0$ is required and the zero mode is treated separately. These
are downstream statistical realizations applied to an already constructed
spectrum.

\subsubsection{Joint compatibility of four readouts}
\label{md:hib:cuatro-lecturas}

Composing the preceding sections under the same hypotheses gives
\begin{equation}
 \boxed{\Xi_\Gamma=\frac{\Omega_\Gamma}{\omega_0}
 =\frac{r_{g,\Gamma}}{R_{108}}
 =\frac{R_{108}}{\bar\lambda_\Gamma}
 =\frac{E_\Gamma}{k_B\Theta_{108}\log3}.}
 \label{md:hib:compatibilidad-cuatro}
\end{equation}
In the circular collection, $\Xi=1$ identifies the section with frequency
$\omega_0$, coincident radii $R_{108}$ and energy $E_{\mathrm{clk}}$.
Its unnormalized thermal weight is $1/3$. This is a compatibility point
of the readouts. The existence of an atlas species with $\Xi=1$, or of a
second generated species for every reciprocal section $1/\Xi$, requires
its generated membership in the atlas; it does not follow from this equality.

\subsection{Energy identity of the ellipse and exact transport}
\label{md:hib:elipse}

For the action ellipse,
\begin{equation}
 a_S=\sqrt{2\hbar}\,e^{\zeta/2},\qquad
 b_S=\sqrt{2\hbar}\,e^{-\zeta/2},\qquad
 Q=a_S\cos\theta,\qquad P=b_S\sin\theta,
 \label{md:hib:elipse-coordenadas}
\end{equation}
with $\zeta=\operatorname{artanh}(C^*_{\deg}/A_{\deg})$, one has
$a_Sb_S=2\hbar$, oriented area
$\mathcal A(\theta)=\hbar\theta$ and $\mathcal A(2\pi)=h$.
The parametric phase $\theta$ and the polar angle of a point on the
ellipse retain their respective geometric meanings.

The declared LC realization uses
\begin{equation}
 L_\zeta=\frac{Z_*}{\omega}e^{-\zeta},\qquad
 C_\zeta=\frac{1}{Z_*\omega}e^\zeta.
 \label{md:hib:lc}
\end{equation}
In action coordinates $Q=\sqrt{Z_*}\,q$ and
$P=\Phi/\sqrt{Z_*}$,
\begin{equation}
 H_0=\frac{\omega}{2}
 \left(e^{-\zeta}Q^2+e^\zeta P^2\right),\qquad
 U_E=\hbar\omega\cos^2\theta,\qquad
 U_B=\hbar\omega\sin^2\theta.
 \label{md:hib:energia-elipse}
\end{equation}
Composition of the two contributions yields $H_0=\hbar\omega$.
The clock and transducer belong to the declared realization; the section
$\hbar$ retains its construction prior to the circuit. The impedance
$Z_{\mathrm{LC}}/Z_*=e^{-\zeta}$ and the constitutive vacuum impedance
$f(s_-)/f(s_+)$ are distinct readouts. Their common angular antecedent does
not establish equality between them.

\subsubsection{Discrete transport of the quadratic form}
\label{md:hib:transporte-discreto}

\begin{proposition}
\label{md:hib:proposicion-discreta}
Let $S_\zeta=\operatorname{diag}(e^{\zeta/2},e^{-\zeta/2})$,
$g_\zeta=S_\zeta^{-T}S_\zeta^{-1}$ and
$\mathsf R(-\theta)$ be a rotation. For shape values $\zeta_n$ and
phases $\theta_n$, the transport
\begin{equation}
 X_{n+1}=S_{\zeta_{n+1}}\mathsf R(-\theta_n)S_{\zeta_n}^{-1}X_n
 =T_nX_n
 \label{md:hib:transporte}
\end{equation}
preserves $I_n=X_n^Tg_{\zeta_n}X_n/2$ exactly and has determinant one.
\end{proposition}
\begin{proof}
Substitution and the identity $\mathsf R^T\mathsf R=I$ give
\begin{equation}
 T_n^Tg_{\zeta_{n+1}}T_n=g_{\zeta_n}.
 \label{md:hib:isometria-transportada}
\end{equation}
Evaluating both sides on $X_n$ shows preservation of the quadratic form.
The factors of $T_n$ have determinant one, whence $\det T_n=1$.
\end{proof}
This composes the previously constructed radial transport. If
$\zeta_n,\theta_n$ are extracted from a TPK route, the identity applies to
that sequence; the algebraic proposition keeps route selection distinct.

\subsubsection{Differentiable transport and connection term}
\label{md:hib:conexion}

If $\zeta(t)$ is of class $C^1$ and $\omega(t)>0$ is continuous, the
relation $X=S_\zeta Y$ and the evolution $\dot Y=-\omega J_0Y$, with
\begin{equation}
 J_0=\begin{pmatrix}0&-1\\1&0\end{pmatrix},
 \label{md:hib:j-canonico}
\end{equation}
give
\begin{equation}
 \dot Q=\frac{\dot\zeta}{2}Q+\omega e^\zeta P,\qquad
 \dot P=-\omega e^{-\zeta}Q-\frac{\dot\zeta}{2}P.
 \label{md:hib:ecuaciones-transporte}
\end{equation}
With the conventions $\dot Q=\partial_PH$ and
$\dot P=-\partial_QH$, the Hamiltonian is, up to a summand depending
only on time,
\begin{equation}
 \boxed{H_{\mathrm{tot}}=
 \frac{\omega}{2}\left(e^{-\zeta}Q^2+e^\zeta P^2\right)
 +\frac{\dot\zeta}{2}QP.}
 \label{md:hib:hamiltoniano-conexion}
\end{equation}
The invariant
\begin{equation}
 I=\frac12\left(e^{-\zeta}Q^2+e^\zeta P^2\right)
 \label{md:hib:invariante}
\end{equation}
is preserved exactly, without requiring an adiabatic approximation.
\begin{proof}
Differentiation gives
\[
 \dot I=\frac{\dot\zeta}{2}
 \left(e^\zeta P^2-e^{-\zeta}Q^2\right)
 +e^{-\zeta}Q\dot Q+e^\zeta P\dot P.
\]
Substitution of \eqref{md:hib:ecuaciones-transporte} cancels the terms
proportional to $\dot\zeta$ and the two $\omega QP$ terms.
Thus $\dot I=0$. If the mixed term in
\eqref{md:hib:hamiltoniano-conexion} is omitted, the remainder is
\begin{equation}
 \dot I=\frac{\dot\zeta}{2}
 \left(e^\zeta P^2-e^{-\zeta}Q^2\right).
 \label{md:hib:defecto-sin-conexion}
\end{equation}
Reversing the sign of the mixed term doubles this defect.
\end{proof}
On the ellipse, $QP=I\sin(2\theta)$ and $\dot\theta=-\omega$, so
\begin{equation}
 H_{\mathrm{tot}}=I\omega+\frac{I\dot\zeta}{2}\sin(2\theta).
 \label{md:hib:hamiltoniano-elipse}
\end{equation}
The matrix of the quadratic form has determinant
$\omega^2-\dot\zeta^2/4$. For $\omega>0$, it is positive definite
precisely when $|\dot\zeta|<2\omega$, and positive semidefinite at equality.
This threshold concerns positivity: existence of the transport and
physical stability retain their respective conditions. Conservation of
$I$ does not imply conservation of $H_{\mathrm{tot}}$, which depends
explicitly on time.

The quantum version requires symmetrization of the mixed term,
$\dot\zeta(QP+PQ)/4$. This expression specifies the required operation
without presupposing an additional quantum representation.

\subsubsection{Compatibility with the action section generating the counterangle}
\label{md:hib:compatibilidad-accion}

In the governing construction,
\begin{equation}
 \zeta=\operatorname{artanh}
 \left[\frac{2(\eta_{\mathrm{ret}}+\alpha)}{1000\alpha}\right].
 \label{md:hib:forma-genealogica}
\end{equation}
Fixing $\alpha$ and $\eta_{\mathrm{ret}}$ also fixes $\zeta$.
The result with variable $\zeta(t)$ is therefore an extension of
transport between shapes. To represent a physical trajectory of the
governing pair, the variation must arise from the corresponding genealogy,
rather than from an independently chosen deformation.

Transport can also be composed exactly when the positive action section
changes between two states while the same dimensional chart is retained.
Let
\begin{equation}
 \hbar_n>0,\qquad
 \rho_{\mathrm{act},n}=\frac{\hbar_{n+1}}{\hbar_n},\qquad
 \widetilde T_n=\sqrt{\rho_{\mathrm{act},n}}\,
 S_{\zeta_{n+1}}\mathsf R(-\theta_n)S_{\zeta_n}^{-1}.
 \label{md:hib:transporte-conforme}
\end{equation}
Then
\begin{equation}
 \det\widetilde T_n=\rho_{\mathrm{act},n},\qquad
 \widetilde T_n^Tg_{\zeta_{n+1}}\widetilde T_n
 =\rho_{\mathrm{act},n}g_{\zeta_n},\qquad
 \frac{I_{n+1}}{\hbar_{n+1}}=\frac{I_n}{\hbar_n}.
 \label{md:hib:accion-normalizada}
\end{equation}
The equalities follow from \eqref{md:hib:isometria-transportada} by
multiplication by the scalar factor. This is conformally symplectic
transport: it preserves normalized action, and it preserves unnormalized
action area precisely when $\rho_{\mathrm{act},n}=1$.
Its differentiable limit has trace $\dot\hbar/\hbar$. A linear Hamiltonian
generator on a plane with a fixed symplectic form has trace zero;
a material variation of $\hbar$ therefore requires specification of the
transported form or the degrees of freedom with which action is exchanged.
The mixed deformation term by itself retains trace zero.

The result is a compatibility criterion for the realizations. To interpret
numerical variations due to a change of units, the bases are transported
first: the chart change preserves the physical section. The composition
does not assert experimental variability of the Planck constant.

\subsubsection{Representation transformation and material deformation}
\label{md:hib:interpretacion-transporte}

If $S_\zeta$ is a passive coordinate change, the mixed term is a
representation connection. The observables are transported as well:
\begin{equation}
 L_X(S_\zeta Y)=L_Y(Y).
 \label{md:hib:observables-transportados}
\end{equation}
This equality preserves the physical readout as the Hamiltonian expression
changes, without introducing an additional interaction or physical work.

A material deformation requires deriving $\zeta$ from the TPK update and
showing a relative change between system and reference, both transported
consistently. One can test a relative phase, a response
$Z_{\mathrm{sistema}}/Z_{\mathrm{referencia}}$ or a transfer between
modes defined by the detector. Two TPK histories inequivalent under a
representation change must be distinguished, and their readout conditions
retained. When an energy exchange is attributed, the balance includes
field, memory and apparatus. The transport theorem delineates this
distinction while keeping separate the dynamical law that realizes it
physically.

\subsection{Momentum, tests and scope of the composition}
\label{md:hib:momento}

The Clifford--Dirac realization receives the route mass. Within the same
action section, its coefficient is written
\begin{equation}
 \frac{m_\Gamma c}{\hbar}
 =\frac{\omega_0\Xi_\Gamma}{c}.
 \label{md:hib:coeficiente-dirac}
\end{equation}
In a free homogeneous realization with the declared relativistic
dispersion relation, the latter is re-expressed as
\begin{equation}
 \Omega^2=c^2|k|^2+\omega_0^2\Xi_\Gamma^2.
 \label{md:hib:dispersion}
\end{equation}
The rewriting preserves the hypotheses of that realization and its
antecedents in the space--time connection. The return moments
$q_-^h-q_+^h$ are tower evaluations and are distinguished from mechanical
momenta.

The factor $\Xi_\Gamma$ reappears as mass weight, relative frequency,
radial dilation and thermal exponent. The constitutive response determines
the common spatial conversion. Barbero retains oriented information
complementary to speed. Action fixes the areal measure; the clock
introduces the temporal rate; shape and its transport specify the
distribution and its evolution.

The corresponding tests comprise mass and frequency ratios; compatibility
of speed, impedance and Barbero; preservation of the quadratic form under
transport; and thermal weights once temperature and multiplicities are
fixed. These relations preserve more data than an isolated decimal
coincidence. Their proofs retain the energy balance and the conditions
of the entropies considered; they do not imply zero thermodynamic entropy.
